% ~750 words
% Five subsections: target audience, gap analysis, adoption ladder, survey, monitoring

The four-group comparison shows that prolific authors exhibit
the same feature profile as occasional contributors
(Section~\ref{sec:four_groups}) --- production volume does not drive
feature exploration.
Instead, the educational context shapes \emph{which} features authors
adopt: Comm-Uni authors gravitate toward narrated, code-heavy content;
Comm-School authors toward multimedia and interactive elements.
The MINT-the-GAP model demonstrates that infrastructure-level choices
can lift adoption for an entire community, but at the cost of narrowing
the explored feature space.
Effective improvements must therefore be \emph{context-sensitive}.

\subsection{Three Adoption Gap Classes}
\label{sec:gap_classes}

The adoption gradient itself can be explained through three gap classes.

\paragraph{Documentation deficit.}
Some features are syntactically non-obvious and lack progressive
examples.
Narrator adoption is high because activation requires a single header
line, whereas animation effects demand a multi-layered secondary syntax
documented only through a comprehensive reference.
The four-group comparison reinforces this: Comm-Uni authors show
significantly lower adoption of animation and script features than
Internal authors, while Comm-School authors --- despite lower narrator
adoption (\SI{60.8}{\percent} vs.\ \SI{88.5}{\percent}) --- achieve
higher rates of multimedia embedding.
The underlying syntax is identical, suggesting that documentation
does not sufficiently bridge the gap for either audience.

\paragraph{Workflow barrier.}
Code projects and executable code require authors to configure external
services and understand multi-file code block interactions ---
a barrier that documentation alone cannot resolve for non-technical
educators.

\paragraph{Configuration lock-in.}
The MINT-the-GAP profile demonstrates that a standardised configuration
--- importing the same libraries (TikZ, Algebrite) across all exercise
files --- can drive near-universal adoption of specific features
(math~\SI{95}{\percent}, imports~\SI{98}{\percent}), but simultaneously
leaves features outside this scope unexplored
(narrator~\SI{9}{\percent}, animations~\SI{0.6}{\percent}).
This pattern represents a design tension: standardised configurations
lower the entry barrier for a predefined feature set but discourage
exploration beyond it.
Future shared configurations should expose optional extension points
so that standardised workflows do not inadvertently narrow the
feature space.

Several features remain below \SI{5}{\percent} in \emph{all} four
groups, including among the core developers themselves
(animated CSS~\SI{2.4}{\percent}, code
projects~\SI{1.3}{\percent}, effects~\SI{0}{\percent}).
Unlike the documentation deficit, where Internal adoption proves that
the feature \emph{can} succeed, these convergent floors suggest that the
features' current design may not match real authoring needs.
They warrant a critical review of whether the feature concepts
themselves require rethinking rather than better documentation.

\subsection{The Adoption Ladder}
\label{sec:ladder}

The four-group analysis reveals that the adoption ladder differs
by educational context.
For \emph{Comm-Uni}, the progression runs:
narrator and links ($>$\,\SI{80}{\percent}) $\to$
code blocks, macros, and tables ($\approx$\,\SI{50}{\percent}) $\to$
animations and imports ($\approx$\,\SI{30}{\percent}) $\to$
embedding features ($<$\,\SI{10}{\percent}).
For \emph{Comm-School}, the ladder is structurally different:
links and images ($>$\,\SI{70}{\percent}) $\to$
macros, narrator, and logo ($\approx$\,\SI{55}{\percent}) $\to$
quizzes, imports, and multimedia ($\approx$\,\SI{35}{\percent}) $\to$
code and advanced embedding ($<$\,\SI{10}{\percent}).
The power-user analysis in Section~\ref{sec:four_groups} shows that
prolific authors exhibit essentially the same feature profile as
occasional contributors, indicating that production volume alone does
not drive feature exploration.
Onboarding documentation should therefore follow these ladders
explicitly and offer context-specific entry points for school vs.\
university authors.

